\chapter{Cycle Space Bond Space}
%%%%%%%%%%%%%%%%%%%%%%%%%%%%%%%%%%%%%%%%%%%%%%%%%%%%%%%%%%%%%%%%%%%%%%%%%%%%%%%
\section{Cycle Space Bond Space}

% Generated by scripts/blueprint_restate.py from TCSlib.GraphTheory.Core.CycleSpaceBondSpace.
% Each body is the STATEMENT only, taken from the declaration's Lean docstring:
% the one-line summary plus the verbatim book statement.  Proof, proof plan,
% significance commentary and formalisation notes are excluded by construction.

\begin{definition}[B\&M's unimodular (Thm 12.3, p. 226): every \emph{full} square submatrix (order $\nu- 1$, selected by an\dots]
\lean{Matrix.IsUnimodular}
\label{Matrix.IsUnimodular}
\leanok
B\&M's \textbf{unimodular} (Thm 12.3, p. 226): every \emph{full} square submatrix (order
\texttt{$\nu$$-$1},
selected by an injective column map) has determinant \texttt{0}, \texttt{+1} or
\texttt{-1}.

\emph{A matrix is said to be unimodular if all
its full square submatrices have determinants \texttt{0}, \texttt{+1} or \texttt{-1}.}
\end{definition}

\begin{definition}[B\&M's oriented incidence matrix $M$ (p. 222)]
\lean{CycleSpace.orientedIncMatrix}
\label{CycleSpace.orientedIncMatrix}
\leanok
B\&M's \textbf{oriented incidence matrix} \texttt{M} (p. 222).

\emph{With each vertex \texttt{v} of \texttt{D} we associate
the function \texttt{m\_v} on \texttt{A} defined by}

m\_v(a) = 1   if a is a link and v is the tail of a
           = -1  if a is a link and v is the head of a
           = 0   otherwise

\emph{The incidence matrix of \texttt{D} is the matrix \texttt{M} whose rows are the
functions \texttt{m\_v}.}
\end{definition}

\begin{definition}[The cycle space $\mathcal{C}$: circulations, i.e]
\lean{CycleSpace.cycleSpace}
\label{CycleSpace.cycleSpace}
\leanok
\uses{CycleSpace.orientedIncMatrix}
The \textbf{cycle space} \texttt{$\mathcal{C}$}: circulations, i.e. \texttt{M *$_{v}$ f
= 0}, B\&M's (12.1).

\emph{Let \texttt{D} be a digraph.  A real-valued
function \texttt{f} on \texttt{A} is called a circulation in \texttt{D} if it satisfies
the conservation
condition at each vertex:}

f$^{-}$(v) = f$^{+}$(v)   for all v $\in$ V                                    (12.1)

\emph{If we think of \texttt{D} as an electrical network, then such a function
\texttt{f} represents a
circulation of currents in \texttt{D}. \dots{} Thus the set of all circulations in
\texttt{D} is a
vector space. We denote this space by \texttt{$\mathcal{C}$}. \dots{} We shall see later
on that each
circulation is a linear combination of the circulations associated with cycles.  For
this reason we refer to \texttt{$\mathcal{C}$} as the cycle space of \texttt{D}.}
\end{definition}

\begin{definition}[The bond space $\mathcal{B}$: potential differences, i.e]
\lean{CycleSpace.bondSpace}
\label{CycleSpace.bondSpace}
\leanok
\uses{CycleSpace.orientedIncMatrix}
The \textbf{bond space} \texttt{$\mathcal{B}$}: potential differences, i.e. \texttt{g =
M$^{T}$ *$_{v}$ p}, B\&M's (12.2).

\emph{Given a function \texttt{p} on the vertex set
\texttt{V} of \texttt{D}, we define the function \texttt{$\delta$p} on the arc set \texttt{A} by the rule that, if an
arc \texttt{a} has tail \texttt{x} and head \texttt{y}, then}

$\delta$p(a) = p(x) - p(y)                                              (12.2)

\emph{If \texttt{D} is thought of as an electrical network with potential \texttt{p(v)}
at \texttt{v}, then,
by (12.2), \texttt{$\delta$p} represents the potential difference along the wires of the
network.
For this reason a function \texttt{g} on \texttt{A} is called a potential difference in
\texttt{D} if
\texttt{g = $\delta$p} for some function \texttt{p} on \texttt{V}.  \dots{}  the set \texttt{$\mathcal{B}$} of all potential differences
in \texttt{D} is a vector space. \dots{} We shall see that each potential difference is
a
linear combination of potential differences associated with bonds.  For this reason
we refer to \texttt{$\mathcal{B}$} as the bond space of \texttt{D}.}
\end{definition}

\begin{definition}[The standard dot-product bilinear form on $K \to K$ (here $A \to F$)]
\lean{CycleSpace.dotForm}
\label{CycleSpace.dotForm}
\leanok
The standard dot-product bilinear form on \texttt{K $\to$ K} (here \texttt{A $\to$ F}).
\end{definition}

\begin{definition}[$G$ is the underlying (simple) graph of the digraph \texttt{(tail, head)}]
\lean{CycleSpace.IsUnderlyingGraph}
\label{CycleSpace.IsUnderlyingGraph}
\leanok
\texttt{G} is the underlying (simple) graph of the digraph \texttt{(tail, head)}.

The \textbf{underlying graph} of a
digraph is obtained by forgetting the direction of every arc.
\end{definition}

\begin{definition}[\texttt{(tail, head)} is an orientation of the simple graph $G$]
\lean{CycleSpace.IsOrientationOf}
\label{CycleSpace.IsOrientationOf}
\leanok
\uses{CycleSpace.IsUnderlyingGraph}
\texttt{(tail, head)} is an orientation of the simple graph \texttt{G}.
\end{definition}

\begin{definition}[An arc-level cycle in the digraph, as a nonempty closed arc-sequence]
\lean{CycleSpace.ArcCycle}
\label{CycleSpace.ArcCycle}
An \textbf{arc-level cycle} in the digraph, as a nonempty closed arc-sequence.
\end{definition}

\begin{definition}[$S \subseteq A$ is acyclic: it contains no arc-cycle]
\lean{CycleSpace.IsArcAcyclic}
\label{CycleSpace.IsArcAcyclic}
\leanok
\uses{CycleSpace.ArcCycle}
\texttt{S $\subseteq$ A} is \textbf{acyclic}: it contains no arc-cycle.
\end{definition}

\begin{definition}[$f_C$: the circulation of an oriented cycle ($\pm 1$ on $C^{+} / C^{-}$, $0$ off $C$)]
\lean{CycleSpace.cycleCirculation}
\label{CycleSpace.cycleCirculation}
\uses{CycleSpace.ArcCycle}
\texttt{f\_C}: the circulation of an oriented cycle (\texttt{$\pm$1} on \texttt{C$^{+}$ / C$^{-}$}, \texttt{0} off \texttt{C}).

\emph{We associate with \texttt{C} the function \texttt{f\_C}
defined by}

f\_C(a) = 1   if a $\in$ C$^{+}$
           = -1  if a $\in$ C \textbackslash{} C$^{+}$
           = 0   if a $\notin$ C

\emph{Clearly, \texttt{f\_C} satisfies (12.1) and hence is a circulation.}
\end{definition}

\begin{definition}[B\&M's edge cut $(S, \bar{S})$: the arcs with exactly one end in $S$]
\lean{CycleSpace.edgeCutSet}
\label{CycleSpace.edgeCutSet}
\leanok
B\&M's edge cut \texttt{[S, $\bar{S}$]}: the arcs with exactly one end in \texttt{S}.

\texttt{[S, $\bar{S}$]} is the set of edges
with one end in \texttt{S} and the other outside.
\end{definition}

\begin{definition}[A bond: a minimal nonempty edge cut]
\lean{CycleSpace.IsBond}
\label{CycleSpace.IsBond}
\leanok
\uses{CycleSpace.edgeCutSet}
A \textbf{bond}: a MINIMAL nonempty edge cut.

\emph{A minimal nonempty edge cut of \texttt{G} is
called a bond.}
\end{definition}

\begin{definition}[$g_B$: the potential difference of a bond, $p =$ indicator of $S$ (B\&M give it, p. 221)]
\lean{CycleSpace.bondPotentialDiff}
\label{CycleSpace.bondPotentialDiff}
\leanok
\texttt{g\_B}: the potential difference of a bond, \texttt{p = } indicator of \texttt{S} (B\&M give it, p. 221).

\emph{Let \texttt{B = [S, $\bar{S}$]} be a bond of \texttt{D}.  We
define \texttt{g\_B} by}

g\_B(a) = 1   if a $\in$ (S, $\bar{S}$)
           = -1  if a $\in$ ($\bar{S}$, S)
           = 0   if a $\notin$ B

\emph{It can be verified that \texttt{g\_B = $\delta$p} where \texttt{p(v) = 1} if
\texttt{v $\in$ S} and \texttt{0} if
\texttt{v $\in$ $\bar{S}$}.}
\end{definition}

\begin{definition}[A basis matrix of a submodule $W \le (A \to F)$: its rows are a basis of $W$]
\lean{CycleSpace.IsBasisMatrix}
\label{CycleSpace.IsBasisMatrix}
\leanok
A \textbf{basis matrix} of a submodule \texttt{W $\le$ (A $\to$ F)}: its rows are a
basis of \texttt{W}.

\emph{A matrix \texttt{B} is called a basis matrix of
\texttt{$\mathcal{B}$} if the rows of \texttt{B} form a basis for \texttt{$\mathcal{B}$}; a basis matrix of \texttt{$\mathcal{C}$} is similarly
defined.}
\end{definition}

\begin{definition}[Restrict Cols]
\lean{CycleSpace.restrictCols}
\label{CycleSpace.restrictCols}
\leanok
\texttt{M | S} --- B\&M's column restriction (\texttt{Matrix.submatrix}).

\emph{If \texttt{R} is a matrix whose columns are labelled with
the elements of \texttt{A}, and if \texttt{S $\subseteq$ A}, we shall denote by
\texttt{R | S} the submatrix of
\texttt{R} consisting of those columns labelled with elements in \texttt{S}.  If \texttt{R} has a
single row, our notation is the same as the usual notation for the restriction of
a function to a subset of its domain.}
\end{definition}

\begin{definition}[A spanning tree of the digraph]
\lean{CycleSpace.IsSpanningTree}
\label{CycleSpace.IsSpanningTree}
A \textbf{spanning tree} of the digraph. Structural predicate; body deferred.

A spanning subgraph that is a tree.
\end{definition}

\begin{definition}[A maximal forest of the digraph]
\lean{CycleSpace.IsMaximalForest}
\label{CycleSpace.IsMaximalForest}
A \textbf{maximal forest} of the digraph. Structural predicate; body deferred.
\end{definition}

\begin{definition}[The fundamental cycle of $a \notin T$ ($T + a$ contains a unique cycle)]
\lean{CycleSpace.fundamentalCycle}
\label{CycleSpace.fundamentalCycle}
\uses{CycleSpace.ArcCycle, CycleSpace.IsMaximalForest}
The \textbf{fundamental cycle} of \texttt{a $\notin$ T} (\texttt{T + a} contains a
unique cycle). path\_unique`.

\emph{Let \texttt{T} be a maximal forest of
\texttt{D}.  Associated with \texttt{T} is a special basis matrix of \texttt{$\mathcal{C}$}.  If \texttt{a} is an arc of
\texttt{$\bar{T}$}, then \texttt{T + a} contains a unique cycle.  Let \texttt{C\_a} denote this cycle and let
\texttt{f\_a} denote the circulation corresponding to \texttt{C\_a}, defined so that \texttt{f\_a(a) = 1}.}
\end{definition}

\begin{definition}[The vertex set $S$ of the fundamental bond of $a \in T$ ($\bar{T} + a$ contains a unique bond, B\&M's\dots]
\lean{CycleSpace.fundamentalBondVertexSet}
\label{CycleSpace.fundamentalBondVertexSet}
The vertex set \texttt{S} of the \textbf{fundamental bond} of \texttt{a $\in$ T}
(\texttt{$\bar{T}$ + a} contains a unique bond,
B\&M's out-of-chapter Thm 2.6).

\emph{Analogously, if \texttt{a} is an arc of
\texttt{T}, then \texttt{$\bar{T}$ + a} contains a unique bond (see theorem 2.6).  Let \texttt{B\_a} denote this
bond and \texttt{g\_a} the potential difference corresponding to \texttt{B\_a}, defined
so that
\texttt{g\_a(a) = 1}.}
\end{definition}

\begin{definition}[$B$ is the tree-$T$ basis matrix of the bond space (rows indexed by $T$)]
\lean{CycleSpace.IsBasisMatrixOfTree}
\label{CycleSpace.IsBasisMatrixOfTree}
\texttt{B} is the tree-\texttt{T} basis matrix of the bond space (rows indexed by \texttt{T}).

\emph{The \texttt{($\nu$-$\omega$) $\times$ $\varepsilon$} matrix \texttt{B} whose
rows are \texttt{g\_a}, \texttt{a $\in$ T}, is a basis matrix of \texttt{$\mathcal{B}$},
called the basis matrix of \texttt{$\mathcal{B}$}
corresponding to \texttt{T}.}
\end{definition}

\begin{definition}[$C$ is the tree-$T$ basis matrix of the cycle space (rows indexed by $\bar{T}$)]
\lean{CycleSpace.IsBasisMatrixOfTree'}
\label{CycleSpace.IsBasisMatrixOfTree'}
\texttt{C} is the tree-\texttt{T} basis matrix of the cycle space (rows indexed by \texttt{$\bar{T}$}).

\emph{The \texttt{($\varepsilon$-$\nu$+$\omega$) $\times$ $\varepsilon$} matrix
\texttt{C}
whose rows are \texttt{f\_a}, \texttt{a $\in$ $\bar{T}$}, is a basis matrix of
\texttt{$\mathcal{C}$}. This follows from the fact
that each row is a circulation and that \texttt{rank C = $\varepsilon$ - $\nu$ +
$\omega$} (because \texttt{C | $\bar{T}$} is an
identity matrix). We refer to \texttt{C} as the basis matrix of \texttt{$\mathcal{C}$}
corresponding to \texttt{T}.}
\end{definition}

\begin{definition}[$\tau(G)$: the number of spanning trees]
\lean{CycleSpace.tau}
\label{CycleSpace.tau}
\uses{CycleSpace.IsSpanningTree}
\texttt{$\tau$(G)}: the number of spanning trees.

\texttt{$\tau$(G)} denotes the number of spanning trees of \texttt{G}.
\end{definition}

\begin{theorem}[The bond space is the row space of the incidence matrix]
\lean{CycleSpace.bondSpace_eq_rowSpace}
\label{CycleSpace.bondSpace_eq_rowSpace}
\leanok
\uses{CycleSpace.bondSpace, CycleSpace.orientedIncMatrix}
Let $D$ be a digraph with finite vertex set $V$ and finite arc set $A$, each arc $a$
having a tail and a head in $V$, and let $F$ be a field. Let $M$ be the oriented
incidence matrix of $D$ over $F$, the $V \times A$ matrix whose row $m_v$ records, for
each arc, whether $v$ is its tail ($+1$), its head ($-1$), or neither ($0$). Then the
bond space $\mathcal{B} \subseteq F^{A}$ of $D$ — the space of potential differences —
coincides with the row space of $M$, that is, with the subspace of $F^{A}$ spanned by
the rows $\{m_v : v \in V\}$.
\end{theorem}

\begin{theorem}[Cycle space as the orthogonal complement of the bond space]
\lean{CycleSpace.cycleSpace_eq_orthogonal_bondSpace}
\label{CycleSpace.cycleSpace_eq_orthogonal_bondSpace}
\leanok
\uses{CycleSpace.bondSpace, CycleSpace.cycleSpace, CycleSpace.dotForm}
Let $D$ be a directed graph with finite vertex set $V$ and finite arc set $A$, and let
$F$ be a field. Equip the space $F^A$ of $F$-valued functions on the arcs with the
standard dot-product bilinear form $\langle f, g\rangle = \sum_{a \in A} f(a)\,g(a)$.
Then the cycle space $\mathcal{C} \subseteq F^A$ (the circulations) is exactly the
orthogonal complement of the bond space $\mathcal{B} \subseteq F^A$ (the potential
differences) with respect to this form; that is, $\mathcal{C} = \{\, f \in F^A : \langle
f, g\rangle = 0 \text{ for all } g \in \mathcal{B} \,\}$.
\end{theorem}

\begin{theorem}[Nonzero circulations contain an arc-cycle]
\lean{CycleSpace.exists_arcCycle_subset_support_of_isCirculation}
\label{CycleSpace.exists_arcCycle_subset_support_of_isCirculation}
\leanok
\uses{CycleSpace.ArcCycle, CycleSpace.cycleSpace}
Let $D$ be a digraph with finite vertex set $V$ and finite arc set $A$, and let $F$ be a
field. If $f \colon A \to F$ is a circulation on $D$ — that is, $f$ lies in the cycle
space $\mathcal{C}$ — and $f$ is not identically zero, then there exists an arc-cycle
$c$ in $D$ (a nonempty closed arc-sequence) all of whose arcs lie in the support of $f$,
i.e. every arc appearing in $c$ satisfies $f(a) \neq 0$.
\end{theorem}

\begin{theorem}[Every nonzero potential difference contains a bond in its support]
\lean{CycleSpace.exists_isBond_subset_support_of_mem_bondSpace}
\label{CycleSpace.exists_isBond_subset_support_of_mem_bondSpace}
\leanok
\uses{CycleSpace.IsBond, CycleSpace.bondSpace}
Let $D$ be a directed graph with finite vertex set $V$ and finite arc set $A$, and let
$F$ be a field. Suppose $g \colon A \to F$ is a potential difference in $D$ — that is,
$g$ lies in the bond space $\mathcal{B}$ — and that $g$ is not identically zero. Then
there is a bond $B$ of $D$ (a minimal nonempty edge cut) all of whose arcs lie in the
support of $g$, i.e. $B \subseteq \{a \in A : g(a) \neq 0\}$.
\end{theorem}

\begin{theorem}[Independent columns of a bond-space basis matrix mark acyclic arc sets]
\lean{CycleSpace.basisMatrix_bondSpace_cols_linearIndependent_iff}
\label{CycleSpace.basisMatrix_bondSpace_cols_linearIndependent_iff}
\uses{CycleSpace.IsArcAcyclic, CycleSpace.IsBasisMatrix, CycleSpace.bondSpace, CycleSpace.restrictCols}
Let $D$ be a digraph with finite vertex set $V$ and finite arc set $A$, each arc $a$
assigned a tail and a head in $V$, and let $F$ be a field. Write $\mathcal{B} \le (A \to
F)$ for the bond space of $D$, the space of potential differences $\delta p$ arising
from functions $p$ on $V$. Suppose $B$ is a basis matrix of $\mathcal{B}$, so that the
rows of $B$ form a basis of $\mathcal{B}$, and let $S \subseteq A$. Then the columns of
the restricted matrix $B \mid S$, indexed by the arcs in $S$, are linearly independent
over $F$ if and only if $S$ is acyclic, i.e. contains no arc-cycle of $D$.
\end{theorem}

\begin{theorem}[Independence of cycle-space basis-matrix columns and bonds]
\lean{CycleSpace.basisMatrix_cycleSpace_cols_linearIndependent_iff}
\label{CycleSpace.basisMatrix_cycleSpace_cols_linearIndependent_iff}
\leanok
\uses{CycleSpace.IsBasisMatrix, CycleSpace.IsBond, CycleSpace.cycleSpace, CycleSpace.restrictCols}
Let $V$ and $A$ be finite sets of vertices and arcs, let each arc $a\in A$ have a tail
$\mathrm{tail}(a)\in V$ and a head $\mathrm{head}(a)\in V$, and let $F$ be a field, so
that the cycle space $\mathcal{C}\le (A\to F)$ is defined. Let $C$ be a basis matrix of
$\mathcal{C}$ — a matrix whose rows form a basis of $\mathcal{C}$ — and let $S\subseteq
A$ be a finite set of arcs. Then the columns of the restriction $C\mid S$ (the submatrix
of $C$ keeping exactly the columns labelled by elements of $S$) are linearly independent
over $F$ if and only if $S$ contains no bond, that is, there is no bond $B$ with
$B\subseteq S$.
\end{theorem}

\begin{theorem}[Dimension of the bond space]
\lean{CycleSpace.finrank_bondSpace}
\label{CycleSpace.finrank_bondSpace}
\leanok
\uses{CycleSpace.IsUnderlyingGraph, CycleSpace.bondSpace}
Let $D$ be a digraph with finite vertex set $V$ and finite arc set $A$, specified by its
tail and head maps $A \to V$, and let $G$ be the underlying simple graph of $D$. Working
over the real field, the bond space $\mathcal{B}$ of $D$ — the space of potential
differences on the arcs — satisfies
\[
\dim \mathcal{B} \;=\; \abs{V} - \omega,
\]
where $\omega$ is the number of connected components of $G$.
\end{theorem}

\begin{theorem}[Dimension of the cycle space]
\lean{CycleSpace.finrank_cycleSpace}
\label{CycleSpace.finrank_cycleSpace}
\leanok
\uses{CycleSpace.IsUnderlyingGraph, CycleSpace.cycleSpace}
Let $D$ be a digraph with finite vertex set $V$ and finite arc set $A$, given by its
tail and head maps, and let $G$ be its underlying simple graph. Then, working over
$\bbr$, the dimension of the cycle space $\mathcal{C}$ of $D$ is
\[
\dim_{\bbr}\mathcal{C} = \abs{A} - \abs{V} + \omega,
\]
where $\omega$ is the number of connected components of $G$.
\end{theorem}

\begin{theorem}[Fundamental circulations of a maximal forest form a basis of the cycle space]
\lean{CycleSpace.isBasisMatrix_cycleSpace_of_maximalForest}
\label{CycleSpace.isBasisMatrix_cycleSpace_of_maximalForest}
\uses{CycleSpace.IsBasisMatrix, CycleSpace.IsMaximalForest, CycleSpace.cycleCirculation, CycleSpace.cycleSpace, CycleSpace.fundamentalCycle}
Let $D$ be a digraph with finite vertex set $V$ and finite arc set $A$, and let $F$ be a
field. Fix a maximal forest $T$ of $D$. For each arc $a \notin T$, the subgraph $T + a$
contains a unique cycle, and we write $f_a$ for the $F$-valued circulation associated
with this fundamental cycle. Then the family $(f_a)_{a \notin T}$, viewed as the rows of
a matrix whose row index ranges over the arcs outside $T$, is a basis matrix of the
cycle space $\mathcal{C}$ of $D$ over $F$: these circulations are linearly independent
and their span is all of $\mathcal{C}$.
\end{theorem}

\begin{theorem}[Basis matrix of the bond space from a maximal forest]
\lean{CycleSpace.isBasisMatrix_bondSpace_of_maximalForest}
\label{CycleSpace.isBasisMatrix_bondSpace_of_maximalForest}
\uses{CycleSpace.IsBasisMatrix, CycleSpace.IsMaximalForest, CycleSpace.bondPotentialDiff, CycleSpace.bondSpace, CycleSpace.fundamentalBondVertexSet}
Let $D$ be a finite digraph with vertex set $V$ and arc set $A$, specified by its tail
and head maps $A \to V$, and let $F$ be a field. Let $T$ be a maximal forest of $D$, and
for each arc $a \in T$ let $g_a$ denote the potential difference of the fundamental bond
of $a$, that is, $g_a = \delta p_a$ where $p_a$ is the indicator function of the vertex
set of that bond. Then the matrix over $F$ whose rows are the functions $g_a$, indexed
by the arcs $a \in T$, is a basis matrix of the bond space $\mathcal{B}$ of $D$;
equivalently, the family $(g_a)_{a \in T}$ is a basis of $\mathcal{B}$.
\end{theorem}

\begin{theorem}[Unimodularity of the tree basis matrix of the bond space]
\lean{CycleSpace.isUnimodular_basisMatrix_bondSpace}
\label{CycleSpace.isUnimodular_basisMatrix_bondSpace}
\uses{CycleSpace.IsBasisMatrixOfTree, CycleSpace.IsSpanningTree, Matrix.IsUnimodular}
Let a finite digraph be given by its tail and head maps on a finite arc set $A$, and let
$T$ be a spanning tree of this digraph. If $B$ is the integer basis matrix of the bond
space corresponding to $T$ — its rows indexed by the arcs of $T$ — then $B$ is
unimodular; that is, every full square submatrix of $B$ (a submatrix whose columns are
picked out by an injective column selection) has determinant $0$, $+1$, or $-1$.
\end{theorem}

\begin{theorem}[Number of spanning trees as $\det(BB^{\mathsf T})$]
\lean{CycleSpace.tau_eq_det_mul_transpose}
\label{CycleSpace.tau_eq_det_mul_transpose}
\uses{CycleSpace.IsBasisMatrixOfTree, CycleSpace.IsSpanningTree, CycleSpace.tau}
Let $D$ be a finite digraph with vertex set $V$ and arc set $A$, specified by its tail
and head maps $A \to V$, and let $T \subseteq A$ be a spanning tree of $D$. Let $B$ be
the integer basis matrix of the bond space corresponding to $T$, with rows indexed by
the arcs of $T$ and columns indexed by all of $A$. Then the number $\tau(D)$ of spanning
trees of $D$ satisfies \[ \tau(D) = \det\!\left(B B^{\mathsf T}\right). \]
\end{theorem}

\begin{theorem}[Unimodularity of the tree basis matrix of the cycle space]
\lean{CycleSpace.isUnimodular_basisMatrix_cycleSpace}
\label{CycleSpace.isUnimodular_basisMatrix_cycleSpace}
\leanok
\uses{CycleSpace.IsBasisMatrixOfTree', CycleSpace.IsSpanningTree, Matrix.IsUnimodular}
Let $T$ be a spanning tree of a finite digraph on vertex set $V$ with arc set $A$, and
let $C$ be the integer basis matrix of the cycle space corresponding to $T$ — the matrix
whose rows are indexed by the arcs outside $T$ and whose columns are indexed by all of
$A$. Then $C$ is unimodular: every full square submatrix of $C$ formed by selecting a
set of columns equal in number to the rows has determinant $0$, $+1$, or $-1$.
\end{theorem}

\begin{theorem}[Spanning-tree count as a Gram determinant of the cycle-space basis matrix]
\lean{CycleSpace.tau_eq_det_mul_transpose_cycleSpace}
\label{CycleSpace.tau_eq_det_mul_transpose_cycleSpace}
\uses{CycleSpace.IsBasisMatrixOfTree', CycleSpace.IsSpanningTree, CycleSpace.tau}
Let $G$ be a directed graph with finite vertex set $V$ and finite arc set $A$, each arc
assigned a tail and a head in $V$, and let $\tau$ denote the number of spanning trees of
$G$. Fix a spanning tree $T \subseteq A$, and let $C$ be the integer matrix over the
cycle space of $G$ that is the basis matrix corresponding to $T$, so that its rows are
indexed by the arcs $a \notin T$ and its columns by the arcs of $A$. Then
\[
\tau = \det\!\left(C\,C^{\mathsf{T}}\right).
\]
\end{theorem}

\begin{theorem}[Spanning-tree count as a determinant of stacked bond- and cycle-space basis matrices]
\lean{CycleSpace.tau_eq_det_fromBlocks}
\label{CycleSpace.tau_eq_det_fromBlocks}
\leanok
\uses{CycleSpace.IsBasisMatrixOfTree, CycleSpace.IsBasisMatrixOfTree', CycleSpace.IsSpanningTree, CycleSpace.tau}
Let $G$ be a finite digraph with vertex set $V$ and arc set $A$, and let $T$ be a
spanning tree of $G$. Let $B$ be the basis matrix of the bond space $\mathcal{B}$
corresponding to $T$, whose rows are indexed by the arcs of $T$, and let $C$ be the
basis matrix of the cycle space $\mathcal{C}$ corresponding to $T$, whose rows are
indexed by the arcs not in $T$; both have integer entries and columns indexed by all of
$A$. Form the square integer matrix $M$ with rows indexed by $A$ by assigning to each
arc of $T$ its row of $B$ and to each remaining arc its row of $C$. Then the number
$\tau(G)$ of spanning trees of $G$ satisfies $\tau(G) = \det M$ or $\tau(G) = -\det M$;
that is, $\tau(G) = \abs{\det M}$.
\end{theorem}

\begin{theorem}[The matrix-tree theorem]
\lean{CycleSpace.matrix_tree_theorem}
\label{CycleSpace.matrix_tree_theorem}
\uses{CycleSpace.orientedIncMatrix, CycleSpace.tau}
\textbf{The matrix-tree theorem} (Kirchhoff, 1847). Let $D$ be a connected digraph with
vertex set $V$ and arc set $A$, and let $M$ be its oriented incidence matrix, the $V
\times A$ matrix whose entry in row $v$ and column $a$ is $1$ if $v$ is the tail of $a$,
is $-1$ if $v$ is the head of $a$, and is $0$ otherwise. Fix any vertex $y \in V$, and
let $K$ be the reduced incidence matrix obtained from $M$ by deleting the row indexed by
$y$. Then the number $\tau(D)$ of spanning trees of $D$ satisfies \[ \tau(D) =
\det\!\left(K K^{\mathsf{T}}\right), \] the equality holding in $\bbz$.
\end{theorem}

\begin{theorem}[Unimodularity of the reduced oriented incidence matrix]
\lean{CycleSpace.isUnimodular_deleteRow}
\label{CycleSpace.isUnimodular_deleteRow}
\leanok
\uses{CycleSpace.orientedIncMatrix, Matrix.IsUnimodular}
Let $D$ be a finite digraph with vertex set $V$ and arc set $A$, where each arc $a$ has
a tail and a head in $V$, and let $M$ be its oriented incidence matrix over $\bbz$: the
rows are indexed by $V$, the columns by $A$, and the $(v,a)$ entry equals $1$ if $v$ is
the tail of $a$, equals $-1$ if $v$ is the head of $a$, and equals $0$ otherwise. Fix a
vertex $y \in V$, and let $K$ be the matrix obtained from $M$ by deleting the row
indexed by $y$. Then $K$ is unimodular: for every injective map $g$ from an index set
into the columns of $K$, the square submatrix of $K$ formed by selecting the columns $g$
(against all remaining rows) has determinant $0$, $1$, or $-1$.
\end{theorem}

\begin{theorem}[Laplacian as oriented incidence matrix times its transpose]
\lean{CycleSpace.lapMatrix_eq_orientedIncMatrix_mul_transpose}
\label{CycleSpace.lapMatrix_eq_orientedIncMatrix_mul_transpose}
\leanok
\uses{CycleSpace.IsOrientationOf, CycleSpace.orientedIncMatrix}
Let $G$ be a simple graph on a finite vertex set $V$, and let $(\mathrm{tail},
\mathrm{head})$ be an orientation of $G$ on a finite arc set $A$, so that each edge of
$G$ is directed from its tail to its head. Write $M$ for the oriented incidence matrix
of this orientation over $\bbz$, the $V \times A$ integer matrix whose entry in row $v$
and column $a$ is $1$ if $v$ is the tail of $a$, $-1$ if $v$ is the head of $a$, and $0$
otherwise. Then the graph Laplacian $L(G)$ of $G$, taken with integer entries, satisfies
\[
L(G) = M M^{\mathsf{T}}.
\]
\end{theorem}

\begin{theorem}[Matrix–tree theorem via a reduced Laplacian determinant]
\lean{CycleSpace.tau_eq_det_lapMatrix_deleteRowCol}
\label{CycleSpace.tau_eq_det_lapMatrix_deleteRowCol}
\leanok
\uses{CycleSpace.IsOrientationOf, CycleSpace.tau}
Let $G$ be a finite simple graph with vertex set $V$, let $(\mathrm{tail},
\mathrm{head})$ be an orientation of $G$, and fix a vertex $y \in V$. Write $L(G)$ for
the integer Laplacian matrix of $G$ (the vertex degree matrix minus the adjacency
matrix), and let $L(G)_{\widehat{y}}$ denote the reduced Laplacian obtained by deleting
the row and the column indexed by $y$, that is, the submatrix of $L(G)$ whose rows and
columns range over the vertices $v \neq y$. Then the number of spanning trees $\tau$
equals the determinant of this reduced Laplacian, \[ \tau =
\det\big(L(G)_{\widehat{y}}\big), \] as an identity of integers, for every choice of the
deleted vertex $y$.
\end{theorem}

\begin{theorem}[Matrix–Tree theorem via a Laplacian cofactor]
\lean{CycleSpace.tau_eq_lapMatrix_cofactor}
\label{CycleSpace.tau_eq_lapMatrix_cofactor}
\leanok
\uses{CycleSpace.IsOrientationOf, CycleSpace.tau}
Let $G$ be a finite simple graph on a vertex set $V$, and let $(\mathrm{tail},
\mathrm{head})$ be an orientation of $G$ with arc set $A$, so that $G$ is the underlying
graph of the digraph $(\mathrm{tail}, \mathrm{head})$. Write $L$ for the graph Laplacian
of $G$ over $\bbz$, the $V \times V$ matrix whose diagonal entry at $v$ is the degree of
$v$ and whose off-diagonal entry at $(u,v)$ is $-1$ when $u$ and $v$ are adjacent and
$0$ otherwise. Fix vertices $i, j \in V$ and a bijection $e$ between the set of vertices
other than $i$ and the set of vertices other than $j$, and let $M$ be the submatrix of
$L$ whose rows are indexed by the vertices $u \neq i$, whose columns are indexed by the
vertices $u \neq i$ as well, and whose $(u, u')$ entry is $L_{u,\, e(u')}$. Then the
number $\tau$ of spanning trees of the digraph $(\mathrm{tail}, \mathrm{head})$
satisfies either $\tau = \det M$ or $\tau = -\det M$.
\end{theorem}

\begin{theorem}[Total unimodularity of the incidence matrix characterizes bipartiteness]
\lean{CycleSpace.incMatrix_isTotallyUnimodular_iff_isBipartite}
\label{CycleSpace.incMatrix_isTotallyUnimodular_iff_isBipartite}
\leanok
Let $G$ be a finite simple graph, and let $M$ be its vertex–edge incidence matrix over
$\bbz$, whose entry in row $v$ and column $e$ is $1$ when $v$ is an endpoint of the edge
$e$ and $0$ otherwise. Then $M$ is totally unimodular — that is, every square submatrix
of $M$ has determinant $-1$, $0$, or $1$ — if and only if $G$ is bipartite.
\end{theorem}

\begin{theorem}[Shank's criterion for the bond–cycle intersection]
\lean{CycleSpace.finrank_inf_pos_iff_dvd_tau}
\label{CycleSpace.finrank_inf_pos_iff_dvd_tau}
\leanok
\uses{CycleSpace.bondSpace, CycleSpace.cycleSpace, CycleSpace.tau}
Let $D$ be a finite digraph with vertex set $V$ and arc set $A$, specified by its tail
and head maps $A \to V$, and let $F$ be a field whose characteristic is a prime $p$.
Writing $\mathcal{B}$ for the bond space and $\mathcal{C}$ for the cycle space of $D$
over $F$, the intersection $\mathcal{B} \cap \mathcal{C}$ has positive dimension over
$F$ if and only if $p$ divides the number $\tau(D)$ of spanning trees of $D$.
\end{theorem}
